\documentclass[10pt,twocolumn]{article}
\usepackage[margin=1in]{geometry}
\usepackage{amsmath,amssymb}
\usepackage{booktabs}
\usepackage{graphicx}
\usepackage{hyperref}
\usepackage{xcolor}
\usepackage{enumitem}
\usepackage{listings}
\usepackage{caption}
\usepackage{subcaption}
\usepackage{multirow}
\usepackage{array}
\usepackage{url}
\usepackage{microtype}
\usepackage{appendix}

\hypersetup{
  colorlinks=true,
  linkcolor=blue!70!black,
  citecolor=green!50!black,
  urlcolor=blue!70!black,
}

\lstset{
  basicstyle=\ttfamily\small,
  breaklines=true,
  frame=single,
  backgroundcolor=\color{gray!10},
}

% ── Title ───────────────────────────────────────────────────────────────────
\title{\textbf{CSE/DSC 234 Project 2: Schema Linking with Supervised Fine-Tuning}\\
       \large End-to-End SLM Pipeline for NL-to-SQL Parsing}

\author{
  Yvonne Wang \\
  \texttt{yvw001@ucsd.edu}
  \and
  Stephanie Xu \\
  \texttt{sxxu@ucsd.edu}
}

\date{June 2026}

% ── Document ────────────────────────────────────────────────────────────────
\begin{document}
\maketitle

% ── Abstract ────────────────────────────────────────────────────────────────
\begin{abstract}
Given a natural-language question and a database schema, schema linking predicts the set of tables and columns the underlying SQL would reference. We present an end-to-end Supervised Fine-Tuning (SFT) pipeline utilizing a Small Language Model (SLM) under a strict 2B parameter constraint. We conducted two rounds of rigorous experimentation. Round 1 leveraged RapidFire AI to optimize schema serialization without data augmentation. Round 2 scaled the dataset through paraphrasing and custom Coverage Extension Data (CED), while sweeping LoRA capacity, context length, and architecture. Our final pipeline, utilizing Qwen3-1.7B, schema formatting with primary/foreign key annotations, and a recall-oriented post-processing algorithm, achieved a leaderboard score of \textbf{0.5134} on the validation set.
\end{abstract}

% ── 1. Introduction ──────────────────────────────────────────────────────────
\section{Introduction}
\label{sec:intro}

Schema linking is a fundamental sub-task in the natural-language-to-SQL (NL-to-SQL) pipeline. Accurate identification of tables and columns minimizes hallucination in downstream SQL generation and limits the search space for large language models. The task requires taking an NL question and a database schema, and emitting a JSON object representing the correct schema links: \texttt{\{"Table": ["col", ...]\}}.

The pipeline is evaluated based on a set-based, case-insensitive $F1$ metric, combining both Table-level and Column-level scores:
\begin{equation}
    \text{Score} = 0.5 \times \text{Table} + 0.5 \times \text{Column}
\end{equation}
where each level computes the macro-averaged $(\text{Precision} + \text{Recall} + \text{F1})/3$ across all questions.

We document a two-phase optimization process. Phase 1 focuses on representation engineering on a small dataset. Phase 2 introduces significant data augmentation and model tuning, concluding with a post-processing algorithm tailored to the evaluation metric's characteristics.

% ── 2. Background ────────────────────────────────────────────────────────────
\section{Background and Constraints}
\label{sec:background}

The project imposes strict infrastructure and architectural limits: the model must be constrained to $\le$ 2B parameters to simulate deployment in edge or low-latency environments. We are prohibited from using external APIs (e.g., OpenAI GPT-4) during inference, forcing reliance on local compute. The total inference budget is set at 15 minutes for approximately 100 questions. 

The baseline dataset consists of only 301 (question, gold-SQL) pairs across 17 diverse SNAILS databases (ranging from national-park biodiversity to SAP business sub-modules). We derive the gold schema links from the underlying SQL queries using a custom \texttt{sql\_to\_schema\_links.py} parsing script.

% ── 3. Phase 1: Representation ───────────────────────────────────────────────
\section{Phase 1: Representation Engineering}
\label{sec:phase1}

\subsection{Motivation}
Our initial goal was to find the optimal prompt structure and schema serialization format cheaply, without introducing the noise of data augmentation. All Round 1 experiments (Table~\ref{tab:r1}) trained LoRA adapters directly on the original 301 examples, orchestrated via \textbf{RapidFire AI} with metrics logged via MLflow.

\subsection{Serialization and Prompt Sweeps}
We experimented with multiple formats:
\begin{itemize}[leftmargin=*,noitemsep]
  \item \textbf{Compact/Typed}: Baseline formats vs. appending specific SQL data types to each column.
  \item \textbf{Schema Sorting}: Sorting tables and columns alphabetically (A$\to$Z) vs. maintaining original database creation order.
  \item \textbf{Pre-filtering}: Using simple keyword heuristics to filter the schema down to a ``Top 10'' candidate list before passing to the LLM.
\end{itemize}

\begin{table}[h]
\centering
\caption{Selected Round 1 Results (Original 301 examples). Sorted by Leaderboard Score.}
\label{tab:r1}
\resizebox{\columnwidth}{!}{%
\begin{tabular}{l l c c c}
\toprule
\textbf{Config} & \textbf{Change} & \textbf{Table} & \textbf{Col} & \textbf{Score} \\
\midrule
Baseline                & Un-optimized                     & 0.011 & 0.037 & 0.0238 \\
Exp3-Col-Filt           & Keep keyword cols only           & 0.005 & 0.030 & 0.0171 \\
Exp3-Q-Hint             & Add ``Key terms''                & 0.185 & 0.076 & 0.1303 \\
Exp1-Typed              & Add column types                 & 0.270 & 0.188 & 0.2291 \\
\textbf{Exp2-Sorted}    & \textbf{Sort A$\to$Z}            & \textbf{0.317} & \textbf{0.168} & \textbf{0.2425} \\
\bottomrule
\end{tabular}%
}
\end{table}

\subsection{Phase 1 Findings}
Schema serialization dominated performance. Introducing column types and sorting elements alphabetically (\textbf{Exp2-Sorted}) yielded the highest score (\textbf{0.2425}). Conversely, aggressive pre-filtering techniques (\texttt{Top-10}, \texttt{Col-Filtered}) severely starved the model's recall, plunging scores to near zero. Crucially, the ceiling of 0.2425 revealed a \emph{coverage problem}: the model could not link tables in the validation set that it had never seen in the 301 training examples.

% ── 4. Phase 2: Augmentation & Tuning ────────────────────────────────────────
\section{Phase 2: Data \& Model Scaling}
\label{sec:phase2}

\subsection{Data Augmentation Streams}
To attack the coverage gap, we built two augmentation streams:
\begin{enumerate}
  \item \textbf{Paraphrase Augmentation (\texttt{aug\_v2}):} Generated $\sim$3,000 examples by paraphrasing original questions into multiple surface forms while retaining identical gold links.
  \item \textbf{Coverage Extension Data (CED-v2):} Synthesized $\sim$4,000 explicit (question, link) pairs designed to cover \emph{every} table and column across the 17 validation schemas—including 229 tables completely absent from the original 301 examples. This directly targeted zero-shot hallucination failures.
\end{enumerate}

\subsection{Architectural Decisions}
Building upon \texttt{Exp2-Sorted}, our final schema representation (\texttt{pkfk}) sorted elements A$\to$Z and appended \texttt{[PK]}/\texttt{[FK]} relational flags to columns (e.g., \texttt{CRASH(CASEID:int[PK], ...)}).

\begin{table}[h]
\centering
\caption{Selected Round 2 Results with Augmentation.}
\label{tab:r2}
\resizebox{\columnwidth}{!}{%
\begin{tabular}{l p{3.5cm} c c c}
\toprule
\textbf{Config} & \textbf{Key Knobs} & \textbf{Table} & \textbf{Col} & \textbf{Score} \\
\midrule
test1-3  & Qwen2.5-1.5B, pkfk, aug\_v2 & 0.554 & 0.329 & 0.4415 \\
test11   & max\_len=2048, $\alpha$=16  & 0.493 & 0.363 & 0.4278 \\
test14-a & QLoRA 4-bit, r=16         & 0.482 & 0.321 & 0.4012 \\
test18-b & Qwen3, pkfk, $\alpha$=16, len=2048 & 0.528 & 0.353 & 0.4405 \\
\textbf{test20-a} & \textbf{pkfk, $\alpha$=32, CED-v2} & \textbf{0.592} & \textbf{0.406} & \textbf{0.4990} \\
test23   & CED-v3 (Calibration)      & 0.522 & 0.403 & 0.4625 \\
\bottomrule
\end{tabular}%
}
\end{table}

\subsection{Phase 2 Findings}
\textbf{Context Length}: Extending \texttt{max\_length} from 1024 to 2048 was critical; 1024 truncated schemas, causing a 40\% empty-prediction rate. 

\textbf{Model Capacity}: A moderate LoRA capacity ($r{=}16, \alpha{=}32$) over 4 epochs proved optimal. Pushing capacity higher ($r{=}64$) caused catastrophic overfitting to the massive new schema space.

\textbf{The Recall vs. Precision Tension}: 
Our final experiments revealed a fundamental tension within the metric. Configs like \texttt{test23} applied stricter calibration data to increase precision. While it succeeded on specific databases (e.g., NYSED), it caused a net score decrease by sacrificing recall across broader databases. Because the grading function features a partial-credit floor for recall but severely punishes missed tables, \textbf{recall-preserving over-prediction mathematically beats strict precision}. Thus, \texttt{test20-a}, which maintained a high-recall posture, remained our strongest base model.

% ── 5. Final Pipeline ────────────────────────────────────────────────────────
\section{Final Pipeline \& Post-Processing}
\label{sec:final}

The final pipeline loads our best \texttt{test20-a} adapter over the \texttt{Qwen3-1.7B} base model. Inference utilizes \texttt{pkfk} formatting, greedy decoding, and automated JSON repair.

To exploit the metric's reward for recall, we applied a \textbf{training-free post-processing algorithm}:
\begin{enumerate}[noitemsep]
  \item \textbf{Column Padding}: For each predicted table, append up to 3 schema columns whose string tokens explicitly match the user query.
  \item \textbf{Table Hallucination}: Add any un-predicted table whose name token-matches the question \emph{if} it contains at least 2 matching columns. 
\end{enumerate}

This logic lifted our final validation score from \textbf{0.4990 $\to$ 0.5134}.

% ── 6. Conclusion ────────────────────────────────────────────────────────────
\section{Conclusion}
\label{sec:conclusion}
We designed an SLM pipeline capable of robust schema linking. Through systematic ablation, we conclude that: (1) Schema serialization format (sorting, PK/FK flags) dictates model comprehension; (2) Data augmentation targeting missing tables is mandatory for generalization; and (3) Under set-based F1 metrics, ensuring high recall through post-processing serves as an effective mechanism to override precision bottlenecks. 

% ── References ───────────────────────────────────────────────────────────────
\bibliographystyle{plain}
\begin{thebibliography}{9}

\bibitem{luoma2025snails}
K.~Luoma and A.~Kumar.
\newblock SNAILS: Schema Naming Assessments for Improved LLM-Based SQL Inference.
\newblock \emph{Proc. ACM SIGMOD}, 2025.

\bibitem{qwen2.5}
Qwen Team.
\newblock Qwen2.5 Technical Report.
\newblock \emph{arXiv:2309.07597}, 2024.

\bibitem{hu2021lora}
E.~J. Hu, Y.~Shen, P.~Wallis, Z.~Allen-Zhu, Y.~Li, S.~Wang, L.~Wang, and W.~Chen.
\newblock LoRA: Low-Rank Adaptation of Large Language Models.
\newblock \emph{ICLR}, 2021.

\end{thebibliography}

% ── Appendix ─────────────────────────────────────────────────────────────────
\appendix

\section{AI Tools Statement}
\label{app:ai_tools}

We used AI coding assistants (Claude / ChatGPT) for: scaffolding SFT training and evaluation scripts (\texttt{main.py}, \texttt{eval.py}), generating paraphrase augmentation data formats, diagnosing the recall-vs-precision regression via per-question error analysis, and drafting this report. \textbf{RapidFire AI} was used in Phase 1 to orchestrate and log prompt experiments. All design decisions, hyperparameter choices, and final result selections were made by the authors. No external inference-time API is used in the submitted pipeline.

\end{document}